\section{Conventions and entropy inputs}\label{sec:preliminaries}

Fix the parameters of Theorem~\ref{thm:main}, and abbreviate
$\Lambda=\Lambda_L$. Unsubscripted norms denote the operator norm for maps
and the Euclidean norm for vectors; $\norm{\cdot}_1$ and
$\norm{\cdot}_2$ denote trace and Hilbert--Schmidt norms. All Hilbert spaces
are finite dimensional. Empty tensor products have dimension one.
Constants denoted by $C,c$ may change between occurrences. They depend only
on $q,J,\Delta$ and explicitly fixed numerical powers, geometric templates,
or tolerances. A polynomial bound means $CL^c$ with this uniform dependence.

Plane distances are sup-norm distances. For a plane set $T$, write
$T^{+r}=\{x:\dist_\infty(x,T)\le r\}$; sampling always means intersection
with the lattice, and physical sampling includes intersection with $\Lambda$.
Graph neighborhoods, when used for Hamiltonian locality, are neighborhoods
in the induced physical lattice and lie in the corresponding ambient
sup-norm enlargement. Distances to the empty set are infinite. Region
complements in entropy formulas refer to the full physical system, unless
another system is specified. Adjacent region letters denote their disjoint
union.

For a normalized state $\omega$ and a region $A$, let $\rho_{\omega,A}$
be its marginal, and set
\begin{align*}
 S_\omega(A)&=-\Tr\rho_{\omega,A}\log\rho_{\omega,A},\\
 I_\omega(A:C)&=S_\omega(A)+S_\omega(C)-S_\omega(AC),\\
 I_\omega(A:C\mid B)&=S_\omega(AB)+S_\omega(BC)
                    -S_\omega(B)-S_\omega(ABC).
\end{align*}
Logarithms are natural and $0\log0=0$. We omit the state subscript when it
is $\Omega$.

\subsection{The imported area law}

The following is the range-one square-domain specialization of the companion
theorem. Stating its uniformity here is useful because later auxiliary
regions are not themselves assigned gapped Hamiltonians.

\begin{proposition}[Entropy area law, {\cite[Theorem~1.1 and Corollary~1.2]{RepoAreaLaw}}]
\label{prop:area-law}
Under \eqref{eq:model}--\eqref{eq:global-gap}, there is
$C_A=C_A(q,J,\Delta)$ such that every $A\subseteq\Lambda$ satisfies
\[
 S_\Omega(A)\le C_A|\partial_\Lambda A|,
 \qquad
 \partial_\Lambda A=\{\{x,y\}\in E(\Lambda):x\in A,
                                      \ y\notin A\}.
\]
The same constant works for every $L$ and every Hamiltonian in the class.
\end{proposition}

The companion theorem allows any finite induced square-lattice domain, with one
bounded Hermitian summand per designated nonempty support of bounded graph
diameter. In \eqref{eq:model} the diameter bound is one. Its only spectral
assumption is the unique full-system ground vector and
\eqref{eq:global-gap}. These are precisely the hypotheses used here.
For a physical sample $X$ of a square of radius $r$, this implies
$S_\Omega(X)\le C(1+r)$, including squares clipped by the outer boundary.

We also use the companion's sublinear collar and amplification propositions.
Their geometric and exponent requirements are checked explicitly in
Proposition~\ref{prop:coarse}; neither is a PEPS approximation statement.
Appendix~\ref{app:companion} records their exact interfaces, including the
positive constraints and two-family construction underlying the companion
argument. The complete companion is \cite{RepoAreaLaw}.

\subsection{Finite-dimensional information facts}

We use strong subadditivity, the chain rule, monotonicity of mutual
information under discarding systems, and, for a pure state on $ABCD$,
the duality
\[
 I(A:C\mid B)=I(A:C\mid D).
\]
These are standard entropy identities and inequalities; see
\cite{LiebRuskai1973} and the preliminary section of \cite{RepoAreaLaw}.
For example $|S(A\mid B)|\le\log\dim\mathcal H_A$, so the entropy cost
of a bounded set of individual sites is bounded by its cardinality times
$\log q$.

For normalized density operators define
$F(\rho,\sigma)=\norm{\sqrt\rho\sqrt\sigma}_1$ and
$D(\rho\Vert\sigma)=\Tr\rho(\log\rho-\log\sigma)$, with the usual
infinite value when the support condition fails.

We recall the relative-entropy estimate through quantum affinity
\cite[Equation~(2.9)]{CarlenLieb2014} and Uhlmann's purification-overlap
theorem \cite[Sections~2 and~5]{Uhlmann1976}. The finite-dimensional proof
below fixes our root-fidelity convention and includes nonminimal purifying
spaces.

\begin{lemma}\label{lem:fidelity}
For normalized finite-dimensional states,
\begin{equation}\label{eq:fidelity-information}
 F(\rho,\sigma)\ge e^{-D(\rho\Vert\sigma)/2},
 \qquad
 D(\rho_{AC}\Vert\rho_A\otimes\rho_C)=I(A:C),
\end{equation}
Moreover, after enlarging the
purifying space if necessary, an isometry on that space realizes overlap
$F(\rho,\sigma)$ between purifications. In particular a
mutual-information bound $\delta$ gives purifications of the joint state
and product state at vector distance at most $\sqrt\delta$, after phase
choice, since $2(1-e^{-\delta/2})\le\delta$.
\end{lemma}
\begin{proof}
Write the eigenvalues and orthonormal eigenvectors of $\rho,\sigma$ as
$(p_i,|i\rangle)$ and $(q_j,|j'\rangle)$. When $D(\rho\Vert\sigma)$ is
finite, $w_{ij}=p_i|\langle i|j'\rangle|^2$ is a probability distribution
on its positive entries, and
\[
 D(\rho\Vert\sigma)=\sum_{i,j}w_{ij}\log\frac{p_i}{q_j}.
\]
Convexity of the exponential gives
\[
 \Tr\sqrt\rho\sqrt\sigma
 =\sum_{i,j}w_{ij}\exp\!\left(-\frac12\log\frac{p_i}{q_j}\right)
 \ge e^{-D(\rho\Vert\sigma)/2}.
\]
Terms with $p_i=0$ contribute zero. The trace norm is at least the
absolute trace, proving the first inequality; the infinite case is
immediate. The relative-entropy identity follows by taking the trace
against $\log(\rho_A\otimes\rho_C)
=\log\rho_A\otimes I+I\otimes\log\rho_C$ on its support.

Canonical purifications have coefficient matrices $\sqrt\rho$ and
$\sqrt\sigma$. Varying a unitary on a common sufficiently large
purifying space makes their overlap range over
$\Tr(\sqrt\rho\sqrt\sigma\,U)$, up to a harmless transpose convention.
Polar decomposition gives the maximum absolute value
$\norm{\sqrt\rho\sqrt\sigma}_1$. Every other purification is obtained
from the canonical one by an isometry on the purifying factor. Enlarge
the common factor to accommodate the entire original purifying space,
extend the isometries to unitaries there, and restrict the maximizing
unitary to the original factor. This proves the stated isometry form,
including unused directions in a nonminimal purification. Finally,
the squared vector distance of unit purifications with positive overlap
$F$ is $2(1-F)$.
\end{proof}

If $\frac12\norm{\rho-\sigma}_1\le\epsilon\le1$, conditional entropy
obeys
\begin{equation}\label{eq:entropy-continuity}
 |S(A\mid B)_\rho-S(A\mid B)_\sigma|
 \le2\epsilon\log\dim\mathcal H_A
 +(1+\epsilon)h\!\left(\frac{\epsilon}{1+\epsilon}\right),
\end{equation}
where $h$ is binary entropy \cite[Lemma~2]{Winter2016}. The dimension is
that of $A$, regardless of the purifying or conditioning space. We may
also use the weaker consequence $C\sqrt\epsilon\log(e\dim\mathcal H_A)$.
All retained physical logarithmic dimensions are $O(L^2)$; hence increasing
a fixed error power absorbs the continuity loss.

Finally, the gap directly implies
\begin{equation}\label{eq:energy-vector}
 \min_\theta\norm{\psi-e^{i\theta}\Omega}^2
 \le\frac{2}{\Delta}\langle\psi,(H-E_0I)\psi\rangle
\end{equation}
for every normalized $\psi$. Indeed the excitation energy bounds
$\Delta(1-|\langle\Omega,\psi\rangle|^2)$, while phase minimization gives
$2(1-|\langle\Omega,\psi\rangle|)$.
